\documentclass[10pt,a4paper]{amsart}
\usepackage{amsmath,amssymb,amsthm,graphicx}
\usepackage{verbatim}
\newenvironment{reference}{\comment}{\endcomment}
%\newenvironment{reference}{}{}
\newenvironment{slogan}{\comment}{\endcomment}
\newenvironment{history}{\comment}{\endcomment}
% For commutative diagrams we use Xy-pic
\usepackage[all]{xy}

% We use 2cell for 2-commutative diagrams.
\xyoption{2cell}
% Theorem environments.
%
\theoremstyle{plain}
\newtheorem{theorem}[subsection]{Theorem}
\newtheorem{proposition}[subsection]{Proposition}
\newtheorem{lemma}[subsection]{Lemma}

\theoremstyle{definition}
\newtheorem{definition}[subsection]{Definition}
\newtheorem{example}[subsection]{Example}
\newtheorem{exercise}[subsection]{Exercise}
\newtheorem{situation}[subsection]{Situation}

\theoremstyle{remark}
\newtheorem{remark}[subsection]{Remark}
\newtheorem{remarks}[subsection]{Remarks}

\numberwithin{equation}{subsection}

% Macros
%
\def\lim{\mathop{\mathrm{lim}}\nolimits}
\def\colim{\mathop{\mathrm{colim}}\nolimits}
\def\Spec{\mathop{\mathrm{Spec}}}
\def\Hom{\mathop{\mathrm{Hom}}\nolimits}
\def\Ext{\mathop{\mathrm{Ext}}\nolimits}
\def\SheafHom{\mathop{\mathcal{H}\!\mathit{om}}\nolimits}
\def\SheafExt{\mathop{\mathcal{E}\!\mathit{xt}}\nolimits}
\def\Sch{\mathit{Sch}}
\def\Mor{\mathop{\mathrm{Mor}}\nolimits}
\def\Ob{\mathop{\mathrm{Ob}}\nolimits}
\def\Sh{\mathop{\mathit{Sh}}\nolimits}
\def\NL{\mathop{N\!L}\nolimits}
\def\CH{\mathop{\mathrm{CH}}\nolimits}
\def\proetale{{pro\text{-}\acute{e}tale}}
\def\etale{{\acute{e}tale}}
\def\QCoh{\mathit{QCoh}}
\def\Ker{\mathop{\mathrm{Ker}}}
\def\Im{\mathop{\mathrm{Im}}}
\def\Coker{\mathop{\mathrm{Coker}}}
\def\Coim{\mathop{\mathrm{Coim}}}

% Boxtimes
%
\DeclareMathSymbol{\boxtimes}{\mathbin}{AMSa}{"02}

%
% Macros for moduli stacks/spaces
%
\def\QCohstack{\mathcal{QC}\!\mathit{oh}}
\def\Cohstack{\mathcal{C}\!\mathit{oh}}
\def\Spacesstack{\mathcal{S}\!\mathit{paces}}
\def\Quotfunctor{\mathrm{Quot}}
\def\Hilbfunctor{\mathrm{Hilb}}
\def\Curvesstack{\mathcal{C}\!\mathit{urves}}
\def\Polarizedstack{\mathcal{P}\!\mathit{olarized}}
\def\Complexesstack{\mathcal{C}\!\mathit{omplexes}}
% \Pic is the operator that assigns to X its picard group, usage \Pic(X)
% \Picardstack_{X/B} denotes the Picard stack of X over B
% \Picardfunctor_{X/B} denotes the Picard functor of X over B
\def\Pic{\mathop{\mathrm{Pic}}\nolimits}
\def\Picardstack{\mathcal{P}\!\mathit{ic}}
\def\Picardfunctor{\mathrm{Pic}}
\def\Deformationcategory{\mathcal{D}\!\mathit{ef}}
\usepackage{cleveref}
\crefname{originallemma}{original lemma}{original lemmas}
\Crefname{originallemma}{Original lemma}{Original lemmas}
\makeatletter
\newlabel{lemma-descent-finite-presentation}{{lemma-descent-finite-presentation}{source}{}{sourcepointer-1}{}}
\newlabel{lemma-descent-finite-presentation@cref}{{[originallemma][0][]lemma-descent-finite-presentation}{[1][0][]source}{}{}{}}
\makeatother
\crefname{originallemma}{original lemma}{original lemmas}
\Crefname{originallemma}{Original lemma}{Original lemmas}
\makeatletter
\newlabel{lemma-descent-flat}{{lemma-descent-flat}{source}{}{sourcepointer-2}{}}
\newlabel{lemma-descent-flat@cref}{{[originallemma][0][]lemma-descent-flat}{[1][0][]source}{}{}{}}
\makeatother
\crefname{originallemma}{original lemma}{original lemmas}
\Crefname{originallemma}{Original lemma}{Original lemmas}
\makeatletter
\newlabel{lemma-flat-finite-presentation-permanence}{{lemma-flat-finite-presentation-permanence}{source}{}{sourcepointer-3}{}}
\newlabel{lemma-flat-finite-presentation-permanence@cref}{{[originallemma][0][]lemma-flat-finite-presentation-permanence}{[1][0][]source}{}{}{}}
\makeatother
\crefname{originallemma}{original lemma}{original lemmas}
\Crefname{originallemma}{Original lemma}{Original lemmas}
\makeatletter
\newlabel{lemma-flat-permanence}{{lemma-flat-permanence}{source}{}{sourcepointer-4}{}}
\newlabel{lemma-flat-permanence@cref}{{[originallemma][0][]lemma-flat-permanence}{[1][0][]source}{}{}{}}
\makeatother
\crefname{originallemma}{original lemma}{original lemmas}
\Crefname{originallemma}{Original lemma}{Original lemmas}
\makeatletter
\newlabel{lemma-gerbe-with-section}{{lemma-gerbe-with-section}{source}{}{sourcepointer-5}{}}
\newlabel{lemma-gerbe-with-section@cref}{{[originallemma][0][]lemma-gerbe-with-section}{[1][0][]source}{}{}{}}
\makeatother
\crefname{originallemma}{original lemma}{original lemmas}
\Crefname{originallemma}{Original lemma}{Original lemmas}
\makeatletter
\newlabel{lemma-local-structure-gerbe}{{lemma-local-structure-gerbe}{source}{}{sourcepointer-6}{}}
\newlabel{lemma-local-structure-gerbe@cref}{{[originallemma][0][]lemma-local-structure-gerbe}{[1][0][]source}{}{}{}}
\makeatother
\crefname{originalreference}{original reference}{original references}
\Crefname{originalreference}{Original reference}{Original references}
\makeatletter
\newlabel{stacks-properties-lemma-descent-surjective}{{stacks-properties-lemma-descent-surjective}{source}{}{sourcepointer-7}{}}
\newlabel{stacks-properties-lemma-descent-surjective@cref}{{[originalreference][0][]stacks-properties-lemma-descent-surjective}{[1][0][]source}{}{}{}}
\makeatother
\crefname{originalreference}{original reference}{original references}
\Crefname{originalreference}{Original reference}{Original references}
\makeatletter
\newlabel{stacks-properties-lemma-surjective-permanence}{{stacks-properties-lemma-surjective-permanence}{source}{}{sourcepointer-8}{}}
\newlabel{stacks-properties-lemma-surjective-permanence@cref}{{[originalreference][0][]stacks-properties-lemma-surjective-permanence}{[1][0][]source}{}{}{}}
\makeatother
\title{The structure morphism of a gerbe is surjective, flat and locally of finite presentation}
\author{The Stacks Project Authors}
\date{stacks-project:stacks-morphisms.tex:local-snapshot-20261012; GNU FDL 1.2 or later}
\begin{document}
\maketitle
\noindent\textit{Editorial note.} All delivered mathematical text is the author's own, copied from the pinned Stacks Project chapter source sources/stacks-algebraic-stacks-morphisms/pinned-original-source.tex; no statement or proof step was rewritten, reordered or removed.\par\medskip
\noindent\textit{Editorial note.} Preamble provenance: the Stacks Project keeps one single shared preamble for every chapter; this card takes it directly from the staged sources/stacks-algebraic-stacks-morphisms/preamble.tex as three verbatim spans (lines 8-12, 14-18 and 154-217), so no derived file was written and nothing was added or redefined. The wrapper class and the base packages are supplied by the assembler.\par\medskip
\noindent\textit{Editorial note.} Bibliography provenance: neither the counted statement nor its proof contains a citation (this lane rescans the whole span before writing the card), so the card carries no bibliography and no entry was reconstructed.\par\medskip
\noindent\textit{Editorial note.} Selection provenance: this card was opened under the file-based control contract (CONTROL.md), with the mechanical pre-card gate run on its statement and proof interval before the card was written.\par\medskip
\noindent\textit{Allowed standard prerequisites (editorial).}\begin{itemize}
\item The descent statements for surjectivity, flatness and local finite presentation. Invoked at the start of the counted proof (source lines 6294-6297). Allows the morphism to be replaced by a flat finitely presented base change.
\item The local structure of a gerbe as a quotient stack over an algebraic space. Invoked inside the counted proof (source lines 6298-6301). Reduces to the case of a quotient stack.
\item The statement that the quotient morphism of a gerbe with a section is surjective flat and locally of finite presentation, and the permanence of these properties. Invoked at the end of the counted proof (source lines 6302-6315). Provides the properties for the quotient and transports them to the original morphism.
\end{itemize}\medskip
\ifdefined\hypertarget\hypertarget{sourcepointer-1}{}\fi
\noindent\textit{Source reference.} Original lemma lemma-descent-finite-presentation: sources/stacks-algebraic-stacks-morphisms/pinned-original-source.tex, source lines 5833--5833. Navigation pointer only: the counted text refers to the Stacks item lemma-descent-finite-presentation, which is an existing labeled result of this same chapter and is not part of this card. The pointer carries the target's own label as its number and its definition line as its span; its content is not reproduced here.\par\medskip
\ifdefined\hypertarget\hypertarget{sourcepointer-2}{}\fi
\noindent\textit{Source reference.} Original lemma lemma-descent-flat: sources/stacks-algebraic-stacks-morphisms/pinned-original-source.tex, source lines 5389--5389. Navigation pointer only: the counted text refers to the Stacks item lemma-descent-flat, which is an existing labeled result of this same chapter and is not part of this card. The pointer carries the target's own label as its number and its definition line as its span; its content is not reproduced here.\par\medskip
\ifdefined\hypertarget\hypertarget{sourcepointer-3}{}\fi
\noindent\textit{Source reference.} Original lemma lemma-flat-finite-presentation-permanence: sources/stacks-algebraic-stacks-morphisms/pinned-original-source.tex, source lines 5853--5853. Navigation pointer only: the counted text refers to the Stacks item lemma-flat-finite-presentation-permanence, which is an existing labeled result of this same chapter and is not part of this card. The pointer carries the target's own label as its number and its definition line as its span; its content is not reproduced here.\par\medskip
\ifdefined\hypertarget\hypertarget{sourcepointer-4}{}\fi
\noindent\textit{Source reference.} Original lemma lemma-flat-permanence: sources/stacks-algebraic-stacks-morphisms/pinned-original-source.tex, source lines 5433--5433. Navigation pointer only: the counted text refers to the Stacks item lemma-flat-permanence, which is an existing labeled result of this same chapter and is not part of this card. The pointer carries the target's own label as its number and its definition line as its span; its content is not reproduced here.\par\medskip
\ifdefined\hypertarget\hypertarget{sourcepointer-5}{}\fi
\noindent\textit{Source reference.} Original lemma lemma-gerbe-with-section: sources/stacks-algebraic-stacks-morphisms/pinned-original-source.tex, source lines 6183--6183. Navigation pointer only: the counted text refers to the Stacks item lemma-gerbe-with-section, which is an existing labeled result of this same chapter and is not part of this card. The pointer carries the target's own label as its number and its definition line as its span; its content is not reproduced here.\par\medskip
\ifdefined\hypertarget\hypertarget{sourcepointer-6}{}\fi
\noindent\textit{Source reference.} Original lemma lemma-local-structure-gerbe: sources/stacks-algebraic-stacks-morphisms/pinned-original-source.tex, source lines 6248--6248. Navigation pointer only: the counted text refers to the Stacks item lemma-local-structure-gerbe, which is an existing labeled result of this same chapter and is not part of this card. The pointer carries the target's own label as its number and its definition line as its span; its content is not reproduced here.\par\medskip
\ifdefined\hypertarget\hypertarget{sourcepointer-7}{}\fi
\noindent\textit{Source reference.} Original reference stacks-properties-lemma-descent-surjective: sources/stacks-algebraic-stacks-morphisms/pinned-original-source.tex, source lines 6296--6296. Navigation pointer only: the counted text refers to the Stacks item stacks-properties-lemma-descent-surjective, which belongs to another chapter of the project and is not part of this card. The pointer carries the target's own label as its number because this chapter's staged snapshot does not contain the other chapter's tag table; its content is not reproduced here.\par\medskip
\ifdefined\hypertarget\hypertarget{sourcepointer-8}{}\fi
\noindent\textit{Source reference.} Original reference stacks-properties-lemma-surjective-permanence: sources/stacks-algebraic-stacks-morphisms/pinned-original-source.tex, source lines 6312--6312. Navigation pointer only: the counted text refers to the Stacks item stacks-properties-lemma-surjective-permanence, which belongs to another chapter of the project and is not part of this card. The pointer carries the target's own label as its number because this chapter's staged snapshot does not contain the other chapter's tag table; its content is not reproduced here.\par\medskip
\begin{lemma}
\label{lemma-gerbe-fppf}
Let $\pi : \mathcal{X} \to \mathcal{Y}$ be a morphism of algebraic stacks.
If $\mathcal{X}$ is a gerbe over $\mathcal{Y}$, then $\pi$ is surjective,
flat, and locally of finite presentation.
\end{lemma}

\begin{proof}
By
Properties of Stacks, Lemma
\ref{stacks-properties-lemma-descent-surjective}
and
Lemmas \ref{lemma-descent-flat} and
\ref{lemma-descent-finite-presentation}
it suffices to prove to the lemma after replacing $\pi$ by a base change
with a surjective, flat, locally finitely presented morphism
$\mathcal{Y}' \to \mathcal{Y}$. By
Lemma \ref{lemma-local-structure-gerbe}
we may assume $\mathcal{Y} = U$ is an algebraic space and
$\mathcal{X} = [U/G]$ over $U$.
Then $U \to [U/G]$ is surjective, flat, and
locally of finite presentation, see
Lemma \ref{lemma-gerbe-with-section}.
This implies that $\pi$ is surjective, flat, and locally
of finite presentation by
Properties of Stacks,
Lemma \ref{stacks-properties-lemma-surjective-permanence}
and
Lemmas \ref{lemma-flat-permanence} and
\ref{lemma-flat-finite-presentation-permanence}.
\end{proof}

\end{document}
